% -----------------------------------------------------------
\newpage
\subsection{Diseño de Base de Datos}

El sistema emplea una base de datos relacional implementada en PostgreSQL y gestionada
mediante Supabase. Esta base representa la capa OLTP del proyecto, encargada de registrar
la operación diaria del restaurante: estructura organizacional, usuarios, mesas físicas,
mesas virtuales, comensales, órdenes, estados de cocina, aportaciones internas y eventos
de auditoría.

El modelo se diseñó bajo criterios de normalización, integridad referencial y trazabilidad
histórica, con el objetivo de preservar la consistencia de los datos transaccionales y
facilitar su posterior extracción hacia el pipeline analítico Medallion. Además, se aplican
políticas de seguridad a nivel de fila (Row Level Security, RLS) para restringir el acceso
a la información según el rol y alcance organizacional de cada usuario.

\subsubsection{Convenciones del Modelo}

La definición del esquema de datos en Bouquet se rige por un conjunto de convenciones
técnicas orientadas a la consistencia y al análisis posterior:

\begin{itemize}
  \item Los identificadores primarios se representan como cadenas con formato UUID v4,
  mapeadas en Prisma como \texttt{text}/\texttt{String} para asegurar portabilidad y
  unicidad global.
  \item El almacenamiento de valores monetarios se realiza en centavos (tipo \texttt{int}),
  eliminando errores de redondeo asociados a tipos de punto flotante.
  \item Las tasas porcentuales se registran en puntos base o \textit{basis points} (\texttt{int}),
  donde $100\,\text{bps} = 1\%$.
  \item En las entidades donde se requiere preservar historial operativo, se incorpora
  eliminación lógica (\textit{soft delete}) mediante el atributo \texttt{archivedAt}.
  \item Las órdenes y liquidaciones almacenan \textit{snapshots} históricos de productos,
  precios y modificadores aplicados al momento de la transacción, garantizando que
  cambios futuros en el catálogo no alteren registros históricos.
  \item La gestión de \texttt{createdAt} y \texttt{updatedAt} se delega al ORM para
  estandarizar marcas temporales de auditoría básica.
\end{itemize}

\subsubsection{Dominios del Modelo}

El esquema se distribuye funcionalmente en ocho dominios lógicos que agrupan las entidades
del sistema:

\begin{enumerate}
  \item \textbf{Catálogos globales:} Entidades de referencia que estandarizan datos maestros.
  \item \textbf{Control de acceso (RBAC):} Gestión de roles, permisos y jerarquía de privilegios.
  \item \textbf{Identidad:} Administración de usuarios colaboradores y perfiles operativos.
  \item \textbf{Organización:} Representación de cadena, zonas y sucursales.
  \item \textbf{Menú y catálogo:} Definición de productos, variantes, modificadores y configuraciones locales por sucursal.
  \item \textbf{Operación de mesa:} Control de mesas físicas, sesiones virtuales y comensales.
  \item \textbf{Órdenes y cocina:} Registro de órdenes, ítems, estados de preparación y eventos de servicio.
  \item \textbf{Liquidación interna y auditoría:} Registro de aportaciones internas, ajustes autorizados, estado de cuenta y trazabilidad de acciones críticas.
\end{enumerate}

\begin{landscape}
\begin{figure}[H]
  \centering
  \includegraphics[width=\linewidth,height=0.9\textheight,keepaspectratio]{erd/rendered/DB-simple.pdf}
  \caption{Modelo de Entidad-Relación Global (Simplificado) --- Bouquet}
  \label{fig:erd-global}
\end{figure}
\end{landscape}

\subsubsection{Relación entre Modelo Transaccional y Modelo Analítico}

El modelo transaccional descrito funciona como fuente primaria para el pipeline analítico.
Las tablas relacionadas con órdenes, mesas, comensales, estados de cocina y liquidaciones
internas son extraídas periódicamente hacia la capa Bronze. Posteriormente, Apache Spark
transforma estos datos hacia capas Silver y Gold para generar indicadores de negocio.

Esta separación permite mantener la base OLTP optimizada para la operación diaria, mientras
que las consultas analíticas se ejecutan sobre estructuras derivadas, reduciendo el impacto
sobre el rendimiento percibido por los usuarios en tiempo real.

\clearpage
